% concatenated from vadic2.tex, chapter1/1motivation.tex, chapter1/2summary.tex, chapter1/3notation.tex, chapter2/backcm1.tex, chapter2/backcm2.tex, chapter2/backcm3.tex, chapter2/backcm4.tex, chapter3/background1.tex, chapter3/grosszagier.tex, chapter4/crysbasis1.tex, chapter4/crysbasis2.tex, chapter4/crysbasis3.tex, chapter4/crysbasis4.tex, chapter5/cmapplication1.tex, chapter5/cmapplication2.tex, chapter5/cmapplication4.tex
\documentclass[12pt,twoside,a4paper]{article}
\usepackage[english]{babel}%Packages
\usepackage[utf8]{inputenc}
\usepackage{amsmath}
\usepackage{amsfonts}
\usepackage[pagebackref]{hyperref}
\usepackage{amsthm}
\usepackage{thmtools}
\usepackage{amssymb, graphics, setspace}
\usepackage{bm}
\usepackage{tikz-cd}
\usepackage{xr}
\usepackage[nameinlink]{cleveref}
\usepackage{url}
\usepackage{xcolor}

\usepackage{enumitem}
\usepackage{authblk}
\usepackage[title]{appendix}
\usepackage{chngcntr}
\usepackage{apptools}
\usepackage{mathtools}
\usepackage{bookmark}
\usepackage{appendix}
\definecolor{crimsonglory}{rgb}{0.75, 0.0, 0.2}
\definecolor{darkpowderblue}{rgb}{0.0, 0.2, 0.6}
\hypersetup{
	colorlinks   = true, %Colours links instead of ugly boxes
	urlcolor     = darkpowderblue, %Colour for external hyperlinks
	linkcolor    = darkpowderblue, %Colour of internal links
	citecolor   = crimsonglory %Colour of citations
}


\newcommand{\subscript}[2]{$#1 _ #2$}%Changing enumeration styles.
\theoremstyle{plain}%Theorem numbering stuff.
\newtheorem{theorem}{Theorem}[section] % reset theorem numbering for each chapter

\newtheorem{definition}[theorem]{Definition} % definition numbers are dependent on theorem numbers
\newtheorem{exmp}[theorem]{Example} % same for example numbers
\newtheorem{conv}[theorem]{Convention} % definition numbers are dependent on theorem numbers
\newtheorem{prop}[theorem]{Proposition}
\newtheorem{cor}[theorem]{Corollary}
\newtheorem{lemma}[theorem]{Lemma}
\newtheorem{claim}[theorem]{Claim}
\newtheorem{remark}[theorem]{Remark}
\newtheorem{remarks}[theorem]{Remarks}
\newtheorem{conj}[theorem]{Conjecture}
\newtheorem{assumption}[theorem]{Assumption}
\newtheorem{para}[theorem]{Paradigm}

\DeclareMathOperator{\card}{card}%Stuff that are supposed to make my life easier.
\DeclareMathOperator{\res}{res}
\DeclareMathOperator{\id}{id}
\DeclareMathOperator{\homm}{Hom}
\DeclareMathOperator{\Res}{Res}
\DeclareMathOperator{\Inf}{Inf}
\DeclareMathOperator{\Cor}{Cor}
\DeclareMathOperator{\ver}{Ver}
\DeclareMathOperator{\End}{End}
\DeclareMathOperator{\GL}{GL}
\DeclareMathOperator{\SL}{SL}
\DeclareMathOperator{\coker}{Coker}
\DeclareMathOperator{\tors}{tors}
\DeclareMathOperator{\diag}{diag}
\DeclareMathOperator{\rank}{rank}
\DeclareMathOperator{\pr}{pr}
\DeclareMathOperator{\tr}{tr}
\DeclareMathOperator{\cl}{cl}
\DeclareMathOperator{\spec}{Spec}
\DeclareMathOperator{\gal}{Gal}
\DeclareMathOperator{\aut}{Aut}
\DeclareMathOperator{\gr}{Gr}
\DeclareMathOperator{\inv}{inv}
\DeclareMathOperator{\ram}{Ram}
\DeclareMathOperator{\Span}{Span}
\DeclareMathOperator{\LT}{LT}
\DeclareMathOperator{\fin}{fin}
\DeclareMathOperator{\atyp}{atyp}
\DeclareMathOperator{\disc}{disc}
\DeclareMathOperator{\zar}{Zar}
\DeclareMathOperator{\chara}{char}
\DeclareMathOperator{\crys}{crys}
\DeclareMathOperator{\isog}{isog}
\DeclareMathOperator{\ord}{ord}
\DeclareMathOperator{\SP}{SP}
\DeclareMathOperator{\adj}{adj}
\DeclareMathOperator{\arch}{arch}
\DeclareMathOperator{\new}{new}
\DeclareMathOperator{\ssing}{ssing}

\newcommand{\im}[1]{Im(#1)}
\newcommand{\Sum}[2]{\displaystyle\sum_{#1}^{#2}}
\newcommand{\Bigcup}[2]{\displaystyle\cup_{#1}^{#2}}
\newcommand{\Bigsum}[2]{\displaystyle\bigoplus_{#1}^{#2}}
\newcommand{\tensor}[2]{\displaystyle\otimes_{#1}^{#2}}
\newcommand{\Z}{\mathbb{Z}}
\newcommand{\N}{\mathbb{N}}
\newcommand{\C}{\mathbb{C}}
\newcommand{\Q}{\mathbb{Q}}
\newcommand{\B}{\mathbb{B}}
\newcommand{\A}{\mathbb{A}}
\newcommand{\R}{\mathbb{R}}
\newcommand{\F}{\mathbb{F}}
\newcommand{\G}{\mathbb{G}}
\newcommand{\fgj}{\hat{\gamma}^{*}_j}
\newcommand{\drfj}[1]{\hat{\gamma}^{*}_{{#1},B_{dR}}}
\newcommand{\CE}{\mathcal{E}}
\newcommand{\CX}{\mathcal{X}}
\newcommand{\ld}{,\ldots,}
\newcommand{\CP}{\mathcal{P}}
\newcommand{\CO}{\mathcal{O}}
\newcommand{\unr}{unr}
\newcommand{\CF}{\mathcal{F}}
\newcommand{\et}{\acute{e}t}
\newcommand{\proet}{pro\acute{e}t}
\newcommand{\HL}{HL}
\newcommand{\V}{\mathbb{V}}
\usepackage[OT2,T1]{fontenc}%for the Sha symbol
\DeclareSymbolFont{cyrletters}{OT2}{wncyr}{m}{n}
\DeclareMathSymbol{\Sha}{\mathalpha}{cyrletters}{"58}
\usepackage{tocloft}%Table of Contents
\renewcommand\cftsecfont{\small}
\renewcommand\cftsecpagefont{\small}
\renewcommand\cftsubsecfont{\small}
\renewcommand\cftsubsecpagefont{\small}

%For matrices with brakcets.
\newcommand\coolover[2]{\mathrlap{\smash{\overbrace{\phantom{%
					\begin{matrix} #2 \end{matrix}}}^{\mbox{$#1$}}}}#2}

\newcommand\coolunder[2]{\mathrlap{\smash{\underbrace{\phantom{%
					\begin{matrix} #2 \end{matrix}}}_{\mbox{$#1$}}}}#2}

\newcommand\coolleftbrace[2]{%
	#1\left\{\vphantom{\begin{matrix} #2 \end{matrix}}\right.}

\newcommand\coolrightbrace[2]{%
	\left.\vphantom{\begin{matrix} #1 \end{matrix}}\right\}#2}

\makeatletter
\newcommand{\address}[1]{\gdef\@address{#1}}
\newcommand{\email}[1]{\gdef\@email{\url{#1}}}
\newcommand{\@endstuff}{\par\vspace{\baselineskip}\noindent\small
	\begin{tabular}{@{}l}\scshape\@address\\\textit{E-mail address:} \@email\end{tabular}}
\AtEndDocument{\@endstuff}
\makeatother

\title{On the $v$-adic values of G-functions II:\\
	\Large Towards Effective Brauer-Siegel}
	\author{Georgios Papas}
\address{Faculty of Mathematics and Computer Science\\
	The Weizmann Institute of Science\\
	234 Herzl Street, Rehovot 76100, Israel, and\\\\	
	
	Institute for Advanced Study\\
	1 Einstein Drive\\
	Princeton, N.J. 08540\\
	U.S.A.}
\email{georgios.papas@weizmann.ac.il, gpapas@ias.edu}

\newcommand{\mathsym}[1]{{}}
\newcommand{\unicode}[1]{{}}
\begin{document}
	\maketitle
	
	\begin{abstract}This is the second paper in a series by the author, centered on the study of values of G-functions associated to a $1$-parameter family of abelian varieties $f:\CX\rightarrow S$ and a point $s_0\in S(K)$ over some number field $K$. 
		
		Here we study the case where $f:\CX\rightarrow S$ is a family of elliptic curves. Extending work of Andr\'e and Beukers, we construct relations among the values of G-functions in this setting at points whose fibers are CM elliptic curves.\end{abstract}
	
%\part{Introduction}
\section{Introduction}
In this installment of our series, we study the case of a $1$-parameter family of elliptic curves and the values of the associated G-functions at points whose fibers are CM elliptic curves. 
 
\subsection{Main Result}

The main geometric picture we consider is that of a family $f:\CE\rightarrow S$ of elliptic curves over a smooth irreducible curve $S$ defined over a number field $K$. We consider also a point $s_0\in S(K)$ together with a uniformizer $x\in K(S)$ at $s_0$. By work of Y. Andr\'e, see \cite{andre1989g}, we may associate to the above picture a family of G-functions ``centered at $s_0$''.

Given a place $v$ of a number field $L$ we denote by $\iota_v:L\hookrightarrow \C_v$ the embedding associated to it. Towards the above direction, our main result, modulo some technical considerations, may be summarized as the following:
\begin{theorem}\label{naiverelations}
	Let $f:\CE\rightarrow S$ and $s_0$ be as above and assume that $\CE_{s_0}$ is CM. Let $Y_G$ be the family of G-functions associated to the pair $(\CE\rightarrow S,s_0)$. 
	
	If $v\in \Sigma_{K(s)}$, a place of $K(s)$, is such that $s$ is $v$-adically close to $s_0$, then there exists a polynomial $R_{s,v}\in\bar{\Q}[\underline{X}]$ for which $\iota_v(R_{s,v}(Y_G(s)))=0$ and $R_{s,v}(Y_G(x))\neq0$ on the functional level. 
	
	Moreover, if $v$ is a place of good ordinary reduction of the fiber $\CE_{s_0}$ then $R_{s,v}$ is independent of such $v$.
\end{theorem}
\begin{remark} By ``$v$-adically close'' we mean that $s$ is within the radius of convergence of the power series in the family $Y_G$ of G-functions associated to our pair. 
	
\end{remark}

The problem of constructing polynomials, such as the $R_{s,v}$ of \Cref{naiverelations}, in the above setting is the most studied of its kind. Indeed, Y. Andr\'e, see \cite{andremots} and F. Beukers, see \cite{beukers}, have already given independently, and with different methods, partial answers to this. 

In more detail, F. Beukers constructed in \cite{beukers}, under assumptions on the morphism $f:\CE \rightarrow S$, polynomials $R_{s,v}$ as in \Cref{naiverelations} when $v$ is an archimedean place of $K(s)$ and when $v$ is finite place with $\displaystyle{v\not|2}$, $3$. On the other hand, Y. Andr\'e in \cite{andremots} constructed polynomials $R_{s,v}$ as in \Cref{naiverelations} under the assumptions that 
\begin{enumerate}
\item $v$ is a place of supersingular reduction of the ``central fiber'' $\CE_{s_0}$ and 
\item $v$ is over some prime $p\in \Q$ which is unramified in the CM field $\End^{0}_{\bar{\Q}}(\CE_{s_0})$.
\end{enumerate}

Following the work of Andr\'e and Beukers, the truly new $R_{s,v}$ in \Cref{naiverelations} are those for places $v$ such that 
\begin{enumerate}
	\item $v$ is over some prime $p\in \Q$ which ramifies in $\End^{0}_{\bar{\Q}}(\CE_{s_0})$, or 
	\item $v$ is a place over which the central fiber $\CE_{s_0}$ has ordinary reduction.
\end{enumerate}

A striking new feature of our construction is that, for $v$ over which $\CE_{s_0}$ has ordinary reduction, the relations we construct are independent of the place $v$. More precisely, we obtain a single polynomial $R_{s,\ord}$, for which $\iota_v(R_{s,\ord}(Y_G(s)))=0$ for all such places for which $s$ and $s_0$ are $v$-adically close. In particular, this polynomial can be constructed from any one such place $v$ and does not depend on the number of such places. Equivalently, the coefficients of $R_{s,\ord}$ are independent of $v$. This phenomenon appears to reflect a previously unnoticed structural property of the G-functions method, and persists in all other cases considered in this series of papers.

\subsection{Height bounds and applications}\label{section:heightboundsintro}

The modern study of the values of G-functions at points of ``special interest'' has its beginnings at the work of E. Bombieri, see \cite{bombg}, and Y. Andr\'e, see \cite{andre1989g}. The ultimate goal of the ``G-functions method'', developed by Y. Andr\'e in \cite{andre1989g}, is to establish height bounds for such points of special interest.

In the geometric setting described in \Cref{naiverelations} the height bounds one gets are of the form \begin{equation}\label{eq:naiveheightboud}
	h(s)<< \deg(\prod_{v}^{} R_{s,v})^{c_0},
\end{equation}where the product runs over all places in $K(s)$ for which $s$ is $v$-adically close to $s_0$. 

These height bounds are connected with Siegel's lower bounds for the class numbers of imaginary quadratic number fields. Indeed, were we to know that the degree of the polynomial on the right hand side of \eqref{eq:naiveheightboud} was bounded polynomially by $[\Q(s):\Q]$, we would have an effective version of Siegel's bounds in \cite{siegelog} upon combining our height bound with results of Masser-W\"ustholz. For more on this we refer the reader to the introduction of \cite{papaseffbrsieg}.

While being unable to obtain such an upper bound for the aforementioned degree, in \Cref{section:applications} we discuss the implications of our relations to the search for an effective version of Siegel's result from the point of view of the G-functions method. In short, we show that the problem may be reduced to finding a bound for the size of the subset $\CP(s,s_0)$ of $\Sigma_{K(s)}$ that consists of places $v$ for which 
\begin{enumerate}
\item $s$ is $v$-adically close to $s_0$ and 
\item $\CE_{s_0}$ has supersingular reduction over $v$ and 
\item $v$ is over some prime $p\in \Q$ which is unramified in $\End^{0}_{\bar{\Q}}(\CE_{s_0})$. 
\end{enumerate}
Alternatively, we would like to have a measure of dependence on $v$ for precisely those $R_{s,v}$ in \Cref{naiverelations} that were constructed by Y. Andr\'e in \cite{andremots}.
\subsection{Outline of the paper}
We start in \Cref{section:background} with some foundational results on the periods of CM elliptic curves, either the archimedean periods coming from the de Rham-Betti comparison isomorphism or the non-archimedean ones coming from the de Rham-crystalline comparison isomorphism. Following this, in \Cref{section:backgroundgfuns} we briefly review the basic framework through which one may associate a family of G-functions to a $1$-parameter family of elliptic curves. 

Next, in \Cref{section:padicprox} we discuss briefly a generalization of \cite{grosszagier}, due to Lauter-Viray \cite{lauterviray}, and record an application relevant to places of proximity over ramified primes of the CM field of endomorphisms of the central fiber. \Cref{section:cmpadicrels} consists of a proof of \Cref{naiverelations} working on a case-by-case basis. More specifically, we describe the $v$-adic relations for $v$ a place of ordinary reduction of the ``central fiber'' $\CE_{s_0}$ of our family of elliptic curves. We also record here the relations of Andr\'e, in the supersingular case, and Beukers, in the archimedean setting. We conclude this section by describing the relations for places over primes ramified in the field $\End^{0}_{\bar{\Q}}(\CE_{s_0})$.

Finally, in \Cref{section:applications}, we discuss the height bounds that we alluded to in \Cref{section:heightboundsintro}. We conclude our exposition by discussing the implications of our height bounds to the search for an effective version of Siegel's result.

\subsection{Notation}
Let $E$ be an elliptic curve over a number field $K$ and $v\in\Sigma_{K}$ a place of $K$. We write $E_v$ for the base change $E\times_{K} K_v$. If $v$ is a finite place of good reduction we will write $\tilde{E}_v$ for the reduction of $E$ modulo $v$. 

Given a family of power series $\mathcal{Y}:=(y_1\ld y_N)\in K[[x]]$ where $K$ is as above, and $v\in \Sigma_{K}$ is some place of $K$, we will write $R_v(y_{j})$ for the $v$-adic radius of convergence of $y_j$. We also set $R_{v}(\mathcal{Y}):=\min R_v(y_j)$. Given $v$ as above, we write $\iota_v:K\hookrightarrow \C_v$ for the associated embedding into $\C_v$, which will stand for either $\C$ or $\C_p$ depending on whether the place $v$ is archimedean or not. Finally, if $y(x)=\Sum{n=0}{\infty}a_n x^n\in K[[x]]$ is a power series as above we will write $\iota_v(y(x)):=\Sum{n=0}{\infty}\iota_v(a_n)x^n$ for the corresponding power series in $\C_v[[x]]$.\\

\textbf{Acknowledgments:} The author thanks Yves Andr\'e for answering some questions about his work, as well as for pointing him towards \cite{grosszagier} that greatly influenced the direction of the exposition here. The author also thanks Chris Daw for his encouragement and many helpful discussions around the G-functions method. The author finally thanks Francesco Saettone for helpful comments that improved the readability of our exposition.

Work on this project started when the author was supported by Michael Temkin's ERC Consolidator Grant 770922 - BirNonArchGeom. Throughout the majority of this work, the author received funding by the European Union (ERC, SharpOS, 101087910), and by the ISRAEL SCIENCE FOUNDATION (grant No. 2067/23). In the final stages of this work, the author was supported by the Minerva Research Foundation Member Fund while in residence at the Institute for Advanced Study for the academic year $2025$-$26$. 

%\part{Background on period matrices}
\section{Period matrices of CM elliptic curves}\label{section:background}

We start in this section with some general facts about the arithmetic of our main objects of study, CM elliptic curves and their periods.

Associated to any elliptic curve $E$ defined over a number field $K$ we will get, for each place $v\in\Sigma_{K}$, a certain period matrix. This will be nothing but the matrix associated to a comparison isomorphism between two cohomology theories, either de Rham and Betti or de Rham and crystalline, and some bases for the respective cohomology groups. In the sequel it will be useful for us to work with particular bases of $H^1_{dR}(E/K)$. We chose to give these the following:
\begin{definition}\label{hodgebasis}Let $E$ be an elliptic curve defined over a number field $K$.  We call an ordered basis $\Gamma_{dR}(E):=\{\omega,\eta\}$ of $H^1_{dR}(E/K)$ a \textbf{Hodge basis} if the following are true:
	\begin{enumerate}
		\item $\omega\neq 0$ is in the first part of the filtration $F^1=e^{*}\Omega_{E/K}\subset H^1_{dR}(E/K)$, and 
		
		\item $\eta$ is such that $\langle \omega,\eta\rangle=1$ with respect to the Riemann bilinear form.
	\end{enumerate}
\end{definition}

\subsection{Canonical Crystalline bases}\label{section:bases}

For the remainder of this section we let  $E/K$ be a CM elliptic curve defined over a number field $K$ and let $v\in \Sigma_{K}$ be a place of $K$. We will work throughout under the further assumptions\footnote{These can be achieved upon base changing $E$ by a finite extension $L/K$. The ``everywhere good reduction'' condition follows from \cite{serretate}. We also note here that the assumption that $\End^{0}_K(E)=\End^0_{\bar{\Q}}(E)$ is equivalent to the assumption that $F\subset K$, see for example the Remark at the beginning of $\S 10$ in \cite{langellfuns}.}that $E$ has everywhere good reduction, that $F:=\End_{\bar{\Q}}^{0}(E)=\End_{K}^{0}(E)$, where $[F:\Q]=2$ is the CM field of endomorphisms of $E$, and that $F\subset K$. 

As per usual notational conventions in $p$-adic Hodge theory we let $K_{v,0}:=W(k_v)[\frac{1}{p}]$, where $k_v$ is the residue field of $K$ at $v$ and $p$ stands for its characteristic. We also denote by $q:=p^f$ the cardinality of $k_v$. From now on, we will also write \begin{equation}
		H^1_v(E):= \begin{cases}H^1_{\crys}(\tilde{E}_v/W(k_v))\otimes K_{v,0}& \text{if }v\in\Sigma_{K,f} \text{ and} \\
H^1_{B}(E_v^{an},\Q)& \text{if } v\in\Sigma_{K,\infty}. 
		\end{cases}
\end{equation}

Over finite places $v$, since $E$ has good reduction at $v$, we have canonical isomorphisms, which we simply refer to as the ``de Rham-crystalline comparison isomorphism'',
\begin{equation}\label{eq:canoniso}
	\rho_v(E):H^1_{dR}(E_v/K_v)\rightarrow H^1_v(E)\otimes_{K_{v,0}} K_v.
\end{equation}These are due to Berthelot-Ogus, see Theorem $2.4$ of \cite{bertogus}. This mirrors the classical ``de Rham-Betti'' comparison of Grothendieck, for $v\in\Sigma_{K,\infty}$ this time, which we will denote again by 
\begin{equation}\label{eq:canonisobetti}
\rho_v(E):H^1_{dR}(E/K)\otimes_K \C \rightarrow H^1_v(E)\otimes_{\Q} \C.
\end{equation}

	Our goal in this subsection is to associate to each $v\in\Sigma_{K,f}$ as above a ``canonically'' chosen symplectic basis of $H^1_v(E)$, denoted $\Gamma_v(E)$ in the sequel. We shall do this by working in cases depending on the reduction of $E$ at the place $v$.
	
	\begin{definition}Let $E/K$ be a CM elliptic curve as above and $v\in\Sigma_{K,f}$ a finite place of $K$. We say that $v$ is an \textbf{ordinary} (resp. \textbf{supersingular}) \textbf{place for }$E$ if the reduction $\tilde{E}_v$ of $E$ at $v$ is an ordinary (resp. supersingular) elliptic curve over the finite field $k_v$.\end{definition}
	
	\subsubsection{Ordinary reduction}\label{section:ordredbases}
	
We start with the case where $v$ is an ordinary place for $E$. Let us also write $F_v\in \End(\tilde{E}_v)$ for the $q$-th power Frobenius endomorphism of the reduction at $v$, where $q=p^f$ is as above. We write $\phi_{v}$ for the endomorphism of the crystalline object $H^1_v(E)$ induced by $F_v$ via pullback. 
%Note that this will coincide with the $f$-th power of the Frobenius $F_v$ above.
	
	Since $\tilde{E}_v$ is ordinary we know that the slopes of $\phi_v$ are $1$ and $0$. In the case of Dieudonn\'e modules this is classical, see for example the discussion on page $98$ of \cite{demazure}. To deduce this about the action of $\phi_v$ on $H^1_{\crys}$ one needs the comparison between $H^1_{\crys}$ and the Dieudonn\'e module of $\tilde{E}_v$ of Berthelot-Breen-Messing, see \cite{berbreenmessing}.
	
	Letting $\lambda_0$, $\lambda_1$ denote the eigenvalues of $\phi_{v}$ we know that these will satisfy $v(\lambda_i)=i$, see for example Manin's theorem on page $98$ of \cite{demazure}. Now we choose $\gamma_v$, resp. $\delta_v$, to be a non-zero eigenvector of $\phi_v$ with eigenvalue $\lambda_1$, resp. $\lambda_0$. These will necessarily be linearly independent, since $\lambda_0\neq  \lambda_1$. Therefore for the Riemann form $\langle\gamma_v,\delta_v\rangle=\mu\neq 0$. Rescaling one of these vectors, say $\delta_v$, by $\mu^{-1}$ we get a basis of $H^1_{v}(E)$ that is symplectic. 
	
	\begin{definition}\label{defcanord}
		We let $\Gamma_v(E):=\{\gamma_v,\delta_v\}$ be the above symplectic basis of $H^1_{\crys}(\tilde{E}_{v}/W(k_v))\otimes K_{v,0}$ and call it a \textbf{canonical basis }of $H^1_v(E)$.
	\end{definition}
	
	The basis $\Gamma_v(E)$ chosen above is obviously not unique, but rather unique up to a scalar. 
Nevertheless, this basis enjoys a certain ``canonicity for morphisms'', due to its construction reflecting the action of Frobenius. In more detail, we have the following:
\begin{lemma}\label{lemmacrysordbasis}Let $\tilde{E}$ and $\tilde{E}'$ be ordinary elliptic curves over $k_v$ and let $f\in \homm(\tilde{E},\tilde{E}')$ be a non-zero homomorphism. Let $\Gamma_v(\tilde{E})=\{\gamma,\delta\}$ and $\Gamma_v(\tilde{E}')=\{\gamma',\delta'\}$ be chosen as above. 
	
	Then, for its pullback $f_{\crys}:H^1_{\crys}(\tilde{E}'/K_{v,0})\rightarrow H^1_{\crys}(\tilde{E}/K_{v,0})$, there exist $\zeta_0$, $\zeta_1\in K_{v,0}$ such that $f_{\crys}(\gamma')=\zeta_1\cdot \gamma$ and $f_{\crys}(\delta')=\zeta_0\cdot \delta$.
\end{lemma}
\begin{proof} The pullback $f_{\crys}$ of $f$ on the level of crystalline cohomology will be an endomorphism in the category of $F$-isocrystals. In other words, we have $f_{\crys}\circ\phi'_{\crys} =\phi_{\crys}\circ f_{\crys}$, where $\phi'_{\crys}$ stands for the $q$-th power Frobenius of $H^1_{\crys}(\tilde{E}'/K_{v,0})$.
	
	 Applying this to $\gamma'$ it is easy to see that we must have that $\phi_{\crys}( f_{\crys}(\gamma'))= \lambda_1\cdot f_{\crys}(\gamma')$. Therefore, since $\phi_{\crys}$ is linear with distinct eigenvalues, we get that there exists $\zeta_1\in K_0$ such that $f_{\crys}(\gamma^{\prime})=\zeta_1\cdot \gamma$. The existence of $\zeta_0$ for $\delta$ follows similarly.
\end{proof}

\subsubsection{Supersingular reduction}\label{section:supersingbases}

Here we assume that $E$ has supersingular reduction at the finite place $v$, which we consider fixed for now. Throughout \Cref{section:supersingbases}, with \Cref{lemmacmperiods} in mind, we assume that the prime $p:=\chara(k_v)$ is unramified in the quadratic field $F=\End^{0}(E)$. 

The description of the basis in this case is due to Y. Andr\'e in \cite{andremots}, see in particular $\S 5$. For a more in depth description we point the interested reader to $\S 5.3.1$ in \cite{andremsj} and Proposition $5.3.3$ and the subsequent remark there. For the reader's convenience, we include here a summary of this construction of Andr\'e.\\

Let us consider the endomorphism algebra $\End^{0}_{k_v}(\tilde{E}_v)$. Since we know that the reduction at $v$ is supersingular, the algebra $\End^{0}_{\bar{k_v}}(\tilde{E}_v)$ will be the unique quaternion algebra $D_{p,\infty}$ over $\Q$ that is ramified only at $p$ and $\infty$. In particular we may find a finite extension $\F_{q'}/\F_q$ for which $\tilde{E}':=\tilde{E}_{v}\times_{\F_q} \spec(\F_{q'})$ is such that $\End^{0}_{\F_{q'}}(\tilde{E}')=D_{p,\infty}$. 

Writing $K'_{v,0}/K_{v,0}$ for the unique unramified extension with residue field $\F_{q'}$ we then get that $D_{p,\infty}$ acts on the crystalline cohomology $H^1_{\crys}(\tilde{E}'/K'_{v,0})=H^1_{v}(E)\otimes_{K_{v,0}} K'_{v,0}$. Tensoring both sides of the canonical isomorphism \eqref{eq:canoniso} of Berthelot-Ogus by the compositum $K'_v:=K'_{v,0}K_v$ we get \begin{equation}
	H^1_{dR}(E/K)\otimes_{K} K'_v\rightarrow (H^1_{v}(E)\otimes_{K_{v,0}} K'_{v,0})\otimes_{K'_{v,0}} K'_{v}.
\end{equation}

Since we assumed that $p$ is unramified in $F$ we know that $F\subset K_{v,0}$. Furthermore, since $F\hookrightarrow D_{p,\infty}$ we know that $D_{p,\infty}\otimes_{\Q} F\simeq M_2(F)$, see for example Proposition $1.2.3$ in \cite{szamuely}. Fixing the embedding $F\hookrightarrow \C_v$ induced from that of $K_v$, as discussed in $5.3.2-5.3.3$ in \cite{andremsj}, this leads to a canonical $F$-subspace $H^1_B(\tilde{E},F):= M_2(F)\cdot \gamma_v\subset H^1_{\crys}(\tilde{E}_v\times_{\bar{\F}_{q}})$ where $\gamma_v\in H^1_{\crys}(\tilde{E}'/K'_{v,0})$ is some eigenvector for the action of $F$ viewed as a sub-algebra of $D_{p,\infty}$ now. 

One may then find a $\delta_v\in H^1_B(\tilde{E},F)$ which is linearly independent to $\gamma_v$ and such that $\langle \gamma_v,\delta_v\rangle=1$ as explained in \cite{andremots}. As in \cite{andremots} we can furthermore pick $\delta_v$ so that it is a complementary eigenvector for the action of $F$. In other words, if say $F=\Q(\sqrt{-d})$ we may assume that $e_{\crys}(\gamma_v)=\sqrt{-d}\cdot \gamma_v$ and $e_{\crys}(\delta_v)=-\sqrt{-d}\cdot \delta_v$, where $e_{\crys}$ denotes the pullback on the level of crystalline cohomology of $e:=\sqrt{-d}\in F$, viewed as an element of the endomorphism algebra of our elliptic curve.

\begin{definition}\label{defcansuper}
	We call the above pair $\Gamma_v(E):=\{\gamma_v,\delta_v\}$ a \textbf{canonical basis} of $H^1_v(E)$.
\end{definition}



\subsection{Period relations}\label{section:cmperiods}

The main period relations, in either the archimedean or $p$-adic setting, among the periods of a CM elliptic curve are summarized in the following:

\begin{lemma}\label{lemmacmperiods} Let $E$, $F$ be as above and let $v\in \Sigma_{K}$ be a place of $K$. Then there exists a Hodge basis $\Gamma_{dR}(E)$ of $H^1_{dR}(E/K)$ such that 
	\begin{enumerate}
		\item if $v\in \Sigma_{K,\infty}$, there exists $\varpi_v\in \C_v$, such that ,with respect to $\Gamma_{dR}(E)$ and a choice of a symplectic basis of $H^1(E_v^{an}, \C)$, the period matrix of $E$ is of the form \begin{equation}\label{eq:cmperarch}
			\begin{pmatrix}\frac{\varpi_v}{2\pi i}&0\\
				0& \varpi_v^{-1}
			\end{pmatrix}.\end{equation}
		
		\item if $v\in \Sigma_{K,f}$ is a finite place of good reduction of $E$ for which $p=\chara(k_v)$ is unramified in $F/\Q$, there exists $\varpi_v\in \C_v$ such that, with respect to $\Gamma_{dR}(E)$ and a choice of a symplectic basis of $H^1_{\crys}(\tilde{E}_v/W(k_v))\otimes \C_v$, the period matrix of $E$ is of the form \begin{equation}\label{eq:cmpernonarch}
			\begin{pmatrix}\varpi_v&0\\
				0& \varpi_v^{-1}
			\end{pmatrix}.
		\end{equation}
	\end{enumerate}
\end{lemma}
\begin{remark}
	We believe that the assertions here are classical and known to experts in either setting. For instance, these are used in passing by Y. Andr\'e in \cite{andremots} in the archimedean case and the case of primes of supersingular reduction of $E$. For a rather detailed proof in the archimedean case we point the interested reader to \cite{papaseffbrsieg}, Lemma $2.9$. We present here a similar argument in the case of finite places.
\end{remark}
\begin{proof}Let $V_{dR}:=H^1_{dR}(E/K)$ and $H^1_v(E):=H^1_{\crys}(\tilde{E}_v/K_{v,0})$, as per our usual notation. Let us also write $F:=\Q(\sqrt{-d})$ with $d\in \Z$ square-free and let $e_{dR}\in \End(V_{dR})$ and $e_{\crys}\in \End(H^1_v(E))$ denote the elements corresponding to the action of $e=\sqrt{-d}\in F$ in the de Rham and crystalline cohomology groups respectively.
	
	We then get for $V_{dR}$ and $H^1_v(E)$ splittings of the form \begin{equation}\label{eq:splitdr}
		V_{dR}=V^{+}_{dR}\oplus V^{-}_{dR}, \text{ and}
	\end{equation}
	\begin{equation}\label{eq:splitstar}
		H^1_v(E):= W^{+}_v\oplus W^{-}_v,
	\end{equation}induced by the action of $F$. Here $V_{dR}^{+}$ (resp. $W^{+}_v$) denotes the subspace on which $e_{dR}$ (resp. $e_{\crys}$) acts as multiplication by $\sqrt{-d}$ and $V_{dR}^{-}$ (resp. $W^{-}_v$) the subspace on which this action is multiplication by $-\sqrt{-d}$. Note here that in order to have this splitting into proper subspaces of $H^1_v(E)$ it is essential that the prime $p$ is unramified in $F/\Q$. Indeed, the assumption $F\leq K$ paired with the fact that $p$ is unramified implies that $F\leq K_{v,0}$, where $K_{v,0}/\Q_p$ is now unramified by definition.	
	
	Let us write $\rho_v:=\rho_v(E)$ for the de Rham-crystalline isomorphism of \cite{bertogus} and let $\omega$, $\eta$ be non-zero vectors in $V^{+}_{dR}$, $V^{-}_{dR}$ respectively, as well as $\gamma$, $\delta$ non-zero vectors in $W^{+}_v$ and $W^{-}_{v}$ respectively. The linear independence of these pairs guarantees that we may furthermore, which we do from now on, assume that they form symplectic bases of $V_{dR}$, and $H^1_v(E)$ respectively. To see this note that the linear independence of, say, $\omega$, $\eta$ guarantees that $\langle\omega,\eta\rangle=\alpha\neq0$, where $\langle\cdot,\cdot\rangle$ denotes the Riemann form on $V_{dR}$, at which point we may replace $\eta$ by $\alpha^{-1}\cdot \eta$.
	
	Let us write $(a_{i,j})$ for the period matrix corresponding to the choice of these bases via $\rho_v$. By compatibility of the action of $F$ on $H^1_{dR}$ and $H^1_{\crys}$ with $\rho_v$ we get that \begin{center}
		$\alpha_{1,1}(\sqrt{-d})\gamma+\alpha_{1,2}(\sqrt{-d})\delta=\rho_v(e_{dr}\cdot \omega )=$
		
		$= e_{\crys}\cdot \rho_v(\omega)=\alpha_{1,1}(\sqrt{-d})\gamma+\alpha_{1,2}(-\sqrt{-d})\delta,$
	\end{center}and similarly for the action of $e$ on $\eta$. Comparing coefficients we see that we must have $\alpha_{i,j}=0$ if $i\neq j$, so that $(\alpha_{i,j})$ is diagonal.
	
	The algebraicity of the construction of the Riemann form and its compatibility with the de Rham-crystalline comparison, see Chapter $5$ in \cite{berbreenmessing}, gives that $\langle \rho_v(\omega), \rho_v(\omega')\rangle =\langle\omega,\omega'\rangle$ for all $\omega$, $\omega'$. Applying this with $\omega$ and $\eta$ shows that $\alpha_{1,1} \alpha_{2,2}=1$ thus concluding the proof.\end{proof}

\begin{remarks}1. The basis $\Gamma_{dR}(E)$ does not depend on the place $v$ and is the same as that chosen in the proof of Lemma $2.9$ in \cite{papaseffbrsieg}, i.e. the archimedean case of part $1$ of the previous lemma.\\
	
	2. The ``canonical'' bases chosen in the discussion preceding \Cref{defcanord}, in the case of ordinary reduction of $E$, resp. the discussion right before \Cref{defcansuper} in the case of supersingular reduction, can be used in the proof of \Cref{lemmacmperiods}, at least up to permutation. 
	
	To see this, in the case of ordinary reduction, note that by \Cref{lemmacrysordbasis} we will have that $\gamma_v$ and $\delta_v$ will be eigenvectors of $e_{\crys}$ in the notation of the proof. Note that if $e_{\crys}(\gamma_v)=-\sqrt{d}\cdot \gamma_v$ then we let $\gamma:=-\delta_v$ and $\delta:=\gamma_v$, where the ``$-$'' is needed to ensure that $\langle \gamma, \delta\rangle =1$.
	
	In the case of supersingular reduction this follows by construction of the basis in \cite{andremots}, see for example $\S5.c$ there for more details.
\end{remarks}

\subsection{Reductions of CM elliptic curves}\label{section:backgroundcmlv}

Let $E$, defined over a number field $K$, be a CM elliptic curve as in \Cref{section:bases}. The description of the type of reduction of $E$ at places of $K$ is classically known here and is due to Deuring. For a proof we point the interested reader to \cite{langellfuns} $\S 13.$, Theorem $12$. 

Phrased under our assumptions in \Cref{section:bases} this is the following:
\begin{theorem}[Deuring's Reduction Theorem]\label{deuringred}
	Let $E/K$ be a CM elliptic curve that has everywhere good reduction and is such that $F:=\End^{0}_K(E)=\End^{0}_{\bar{\Q}}(E)$. Let $v\in\Sigma_{K,f}$ be a finite place of $K$ and let $p\in \Q$ be the prime over which $v$ lies. 
	
The reduction $\tilde{E}_v$ of $E$ at $v$ is supersingular if and only if $p$ either ramifies or is inert in $F$.\end{theorem}


   % \part{Background on G-functions}
\section{Background on the G-functions method}\label{section:backgroundgfuns}

Here we summarize the connection between G-functions and a $1$-parameter family of elliptic schemes. For a more detailed exposition, we point the interested reader to $\S 3$ of \cite{papaspadicpart1}, whose terminology we have chosen to follow, as well as \cite{daworrpap}.

The geometric picture permeating the rest of the paper is the one outlined in \Cref{naiverelations}. Namely, we consider some $1$-parameter family of elliptic curves $f:\CE\rightarrow S$ defined over a number field $K$, together with a distinguished point $s_0\in S(K)$ whose fiber is a CM elliptic curve, which we denote by $E_0$ for ease of notation. From the perspective of the G-functions method it is convenient to make assumptions about both the ``central fiber'' $E_0$ and the curve $S$, as long as those may be attained by base changing our original picture by either a finite cover $S'$ of $S$ or by tensoring the entire picture by a finite extension $K'/K$.

With this in mind, we assume from now on that $E_0$ has everywhere good reduction, see \cite{serretate}, and that all endomorphisms of $E_0$ are defined over the base field $K$, i.e. that $\End^{0}_{\bar{\Q}}(E_0)=\End^{0}_{K}(E_0)$ is some CM field. 

In this context, of allowing base changes to the original picture, one may find a ``local parameter'' $x\in K(S)$ at $s_0$ and a tuple of power series $\mathcal{Y}\in \bar{\Q}[[x]]^n$ that one can think of as ``power series on the curve $S$ centered at $s_0$''. These power series appear, in a canonical way, as the entries of matricial solution to the system of differential equations one gets by considering the differential module $(H^1_{dR}(\CE/S),\nabla)$ together with a basis of sections of this module on an affine neighborhood of $s_0$. 

In our setting $\mathcal{Y}$, often denoted $Y_G(x)$ in the sequel, is nothing but a collection of matrices in $\GL_2(\bar{\Q}[[x]])$. In practice, the number of such matrices depends on the preimages of our original $s_0$ under a potentially necessary finite cover $S'\rightarrow S$, see \cite{daworr4} where these ideas were first introduced. For reasons of expositional simplicity we have chosen to consider the case where there is only one such matrix in our family, especially since this is the crucial step in the method we employ. 

Considering the values of $Y_G(x)$ at points $s\in S(\bar{\Q})$ makes sense only after we have introduced some topology to our curve, which we do via considering the analytification of our curve $S$, either in the complex analytic or $p$-adic sense, with respect to some place $v$ of $K$. The $v$-adic values $\iota_v(Y_G(x(s)))$ then make sense as long as $|x(s)|_v$ is smaller than the $v$-adic radius of convergence of the family $Y_G(x)$. We thus naively think of $v$-adic analytic discs $\Delta_v$ centered at our $s_0$ on the curve with sufficiently small radii as part of setting.

For finite places $v$, after possibly altering the family $Y_G$, we may select $\Delta_v$ so that if $s\in S(\bar{\Q})\cap \Delta_v$ then $\CE_s$ and $E_0$ have the same reduction modulo $v$. For more on this, see $\S 3.3.2$ of \cite{papaspadicpart1}. With this further property for the $\Delta_v$ in mind for finite places, from now on we adopt the following:
\begin{definition}\label{defnvadicprox}
	Let $s\in S(\bar{\Q})$ and $w\in\Sigma_{K(s)}$ and $v\in\Sigma_{K}$ such that $w|v$. We say that $s$ is $w$-adically close to $s_0$ if $s\in \Delta_v$ for the analytic discs chosen above.
\end{definition}

On such $\Delta_v$ the values of our G-functions attain cohomological significance. In the archimedean setting this is due to work of Andr\'e in \cite{andre1989g} and in the $p$-adic setting the first observation of this phenomenon appears in Andr\'e's \cite{andremots}. In short, we have in our disposal relative versions of the comparison isomorphisms that appear in \Cref{section:background}. In the archimedean setting this is classically due to Grothendieck and in the $p$-adic setting this is due to work of Berthelot-Ogus, see \cite{bertogus}, and Ogus, see \cite{ogus}, in the ramified case. All in all, for $s\in\Delta_v\cap S(\bar{\Q})$ we will have \begin{equation}
	\CP_v(s)=\iota_v(Y_G(x(s)))\CP_v(s_0),
\end{equation}where $\CP_v$ stand for the matrices associated to the aforementioned isomorphism, the choice of a basis of sections of $H^1_{dR}(\CE/S)$, and a basis of the fiber $H^1_v(E_0)$, which we may identify with the horizontal sections of $(H^1_{dR}(\CE/S),\nabla)$ on the disc $\Delta_v$. For more details on this we point the interested reader to the proof of Theorem $3.4$ in \cite{papaspadicpart1}.

In constructing relations among the $v$-adic values of the $Y_G(x)$ at points of interest, we will want to make certain that these are ``non-trivial'' in the terminology of $VII.5$ in \cite{andre1989g}. Our main tool in this direction is the description of the ``trivial relations'' among the members of the family $Y_G(x)$. We record these in the following:

\begin{prop}\label{trivialrelations}
	Let $Y_G(x)$ be the family of G-functions associated to $(\CE\rightarrow S, s_0)$ as above. Then $Y_G(x)^{\zar}\subset (M_{2\times 2})_{\bar{\Q}(x)}$ is the subvariety cut out by the ideal \begin{equation}\label{eq:trivialideal}
		I(\SL_2)=\{\det(X_{i,j})-1 \},
	\end{equation}where $X_{i,j}$ are such that $M_{2\times 2}$ is given by $\spec({\Q}[X_{i,j}:1\leq i,j\leq 2])$.
\end{prop}
This result seems to be classical, it is mentioned without proof in \cite{beukers} for example. For a proof we point the interested reader to $\S 4.5.6$ of \cite{daworrpap}.

\section{$p$-adic proximity}\label{section:padicprox}

The issue of ``$p$-adic proximity'' of $s$ to $s_0$ may be linked directly to the work of Gross-Zagier in \cite{grosszagier}. We record here, for the convenience of the reader, a generalization of their results due to Lauter-Viray. As a consequence of this we conclude with a short lemma we will need in the sequel around $p$-adic proximity of CM points to our center $s_0$, for primes $p$ ramified in the CM field giving the algebra of endomorphisms of the ``central fiber'' $E_{0}$ in the notation of \Cref{section:backgroundgfuns}.

\subsection{A formula of Gross-Zagier}\label{section:lauterviraysummary}
Let $d_j$, for $j=1$, $2$, be the discriminants of two distinct quadratic imaginary orders $\CO_j$ and consider the product 
\begin{equation}
	J(d_1,d_2):=\underset{\underset{\disc(\tau_j)=d_j}{[\tau_1],[\tau_2]}}{\prod}(j(\tau_1)-j(\tau_2))^{\frac{8}{w_1w_2}},
\end{equation}with $[\tau_j]$ ranging through all elements of $\mathcal{H}/\SL_2(\Z)$ with prescribed discriminants and where $w_j$ denote the number of roots of unity in the order $\CO_j$.

In \cite{grosszagier}, Gross and Zagier established a remarkable formula for the quantity $J(d_1,d_2)$ under the assumptions that \begin{enumerate}
	\item $d_1$ and $d_2$ are relatively prime and
	\item the $d_j$ are fundamental discriminants, i.e. discriminants of maximal orders.
\end{enumerate}   This formula of Gross-Zagier was further generalized, fairly recently by Lauter-Viray, see \cite{lauterviray} Theorem $1.1$, to the case where $d_1\neq d_2$ and without any assumptions on the maximality of the corresponding orders. 

As a corollary of their formula, see \cite{grosszagier} Corollary $1.6$, Gross-Zagier get an elegant description of the primes dividing $J(d_1,d_2)$, albeit under their aforementioned assumptions on the pair $(d_1,d_2)$. The natural generalization of this, again due to Lauter-Viray, is the following:
\begin{theorem}[Lauter-Viray, \cite{lauterviray} Corollary $1.3$]\label{lauterviraytheorem}
	Let $d_1\neq d_2$ be two distinct discriminants as above and let $l\in \N$ be a prime with $l|J(d_1,d_2)$. Then\begin{enumerate}
		\item $d_2=d_1\cdot l^{2k}$ for some $k\in \Z\backslash \{0\}$, or
		\item there exists an integer $m=\frac{d_1d_2-x^2}{4}>0$ with $l|m$ such that the Hilbert symbol $(d_j,-m)_p\neq 1$ if and only if $p=l$.
	\end{enumerate}
	
	In particular, if $l\not| d_1d_2$ then $\left( \frac{d_1}{l}\right)=\left( \frac{d_2}{l}\right)=-1$.
\end{theorem}

\subsection{Application to primes of proximity}\label{section:applicationsofgrosszagier}

Here we consider the geometric picture sketched in \Cref{section:backgroundgfuns}. Namely we consider a fixed $1$-parameter family $f:\CE\rightarrow S$ of elliptic curves defined over a number field $K$. As in \Cref{section:backgroundgfuns} we assume that we are also given a point $s_0\in S(K)$ for which $E_{0}:=\CE_{s_0}$ is a CM elliptic curve with everywhere good reduction and with $F_0:=\End^{0}_{\bar{\Q}}(E_0)=\End^{0}_K(E_0)$. 

In the above picture, following \Cref{defnvadicprox} and the conventions around it, we have:
\begin{lemma}\label{lemproximityramified}
	Let $s\in S(\bar{\Q})$ be such that $\CE_s$ is a CM elliptic curve and write $d_P=\disc(\End_{\bar{\Q}}(\CE_P))$, for $P\in \{s_0,s\}$, to be the corresponding discriminant. Let us also fix a finite place $v\in \Sigma_{K,f}$ over a rational prime $p$, where $p$ is ramified in the CM field $F_0=\End^{0}(E_0)$. 
	
	If $s$ is $v$-adically close to $s_0$ then $\End^{0}_{\bar{\Q}}(\CE_s)=F_0$, i.e. $\CE_s$ has CM by the same field as $E_0$.
\end{lemma}
\begin{proof}
	Assume that it is not so. In particular, we must have that $d_s\neq d_{s_0}$. 
	
	By our conventions on the disks $\Delta_v$ we get that $\CE_s$ and $E_0$ have the same reduction over a place $w\in\Sigma_{K(s),f}$ with $w|v$. In particular, $w(j(\CE_s)-j(E_0))>0$ and hence $p|J(d_s,d_{s_0})$. Note that since $p$ is ramified in $F_0$, we must have $p|d_{s_0}$. Now we apply \Cref{lauterviraytheorem} to conclude that $d_s=p^{2k}d_{s_0}$ for some $k\in \Z$ and our result follows.
\end{proof}
% \part{CM elliptic curves}
\section{Relations among values of G-functions}\label{section:cmpadicrels}

This constitutes the main technical part of our text. In short, we describe the construction of relations among the values of G-functions associated to a $1$-parameter family $f:\CE\rightarrow S$ of elliptic curves at points $s\in S(\bar{\Q})$ whose corresponding fibers are CM elliptic curves. 

In essence, this section is a proof of \Cref{naiverelations}. The construction of the relations described here depends on the ``type'' of the place $v$ with respect to the central fiber $E_0$.

\subsection{Notation}\label{section:finrelsnotation}Throughout this section we work in a somewhat simplified form of the general setting described in the previous section. For the remainder of this section we consider fixed a $1$-parameter family of elliptic curves $f:\CE\rightarrow S$, where $S$ is a smooth and geometrically irreducible curve over a number field $K$. We also consider a fixed point $s_0\in S(K)$ for which we assume that the fiber is a CM elliptic curve $E_0:=\CE_{s_0}$.

In order to simplify our exposition, throughout this section we assume that the curve $E_0$ satisfies the following:\begin{enumerate}
	\item $E_0$ has everywhere\footnote{A fact we may assume up to base changing our original morphism $\pi$ by some finite extension of $K$ due to \cite{serretate}.} good reduction,
	
	\item $\End_{\bar{\Q}}(E_0)=\End_K(E_0)=F_0$, and 
	
	\item there exists a uniformizer $x\in K(S)$ that has a simple zero at $s_0$.
\end{enumerate}

For the remainder of this subsection we also fix a place $v\in \Sigma_{K}$ as well as a point $s\in S(K)$ which is $v$-adically close to $s_0$ in the sense discussed in \Cref{section:backgroundgfuns}.

Given the above data, we choose a symplectic basis $\{\omega,\eta\}$ of $H^1_{dR}(\CE/S)$ whose fiber at $s_0$ consists of $F_0$-eigenvectors, as in the proof of \Cref{lemmacmperiods}. We also may choose $\omega$ so that it is a section of $F^1H^1_{dR}(\CE/S)$. This does not interfere with the assumption that the fiber $\omega_{0}$ is an eigenvector for the action of $F_0$, since $F^1H^1_{dR}(E_0/K)$ are precisely the invariant differentials. From \Cref{section:backgroundgfuns} we therefore get a matrix $Y\in \SL_2(\bar{\Q}[[x]])$ of G-functions associated with this choice of basis.


	

\subsection{Relations at ordinary places}\label{section:cmordpadic}
We start with the case where the place $v$ is a finite place of ordinary reduction for the ``central fiber'' $E_{0}$. In this case we have:

\begin{prop}\label{thmordrels}With notation as in \Cref{section:finrelsnotation}, let $Y_G=(Y_{i,j}(x))\in \SL_2(\bar{\Q}[[x]])$ be the matrix of G-functions associated to the pair $(f,s_0)$ as in \Cref{section:backgroundgfuns}. 
	
	Then for all $s\in S(\bar{\Q})$ for which $\End^{0}_{\bar{\Q}}(\CE_s)=\End^{0}_{\bar{\Q}}(E_0)$ and for all finite places of $\Sigma_{K(s)}$ that are ordinary for $E_0$ and are such that $s$ is $v$-adically close to $s_0$ we have $\iota_v(Y_{1,2}(x(s)))=0$.\end{prop}
\begin{proof}
	We begin by noting that our assumption that $E_0$ has ordinary reduction at $v$ allows us to use \Cref{lemmacmperiods}. Indeed by Deuring's criterion, see Theorem $12$ in $\S13.4$ of \cite{langellfuns}, having ordinary reduction implies that the prime $p=\chara{k_v}$ is unramified in $F_0$.
	
	We choose a basis $\{\gamma,\delta\}$ of $H^1_{\crys}(\tilde{E}_{0,v}/K_{0,v})$ as in \Cref{lemmacrysordbasis}. Then, again by \Cref{lemmacmperiods}, since both this as well as the chosen basis $\{\omega_0,\eta_0\}$ of $H^1_{dR}(E_0/K)$ are symplectic and consist of $F_0$-eigenvectors in the respective spaces, the matrix of relative periods associated to the comparison isomorphism $\CP_v$ will be of the form $\CP_v(s)=\iota_v(Y(x(s)))\cdot \begin{pmatrix}
		\varpi_v&0\\0&\varpi_v^{-1}
	\end{pmatrix}$, where $\varpi_v\in \C_v$.

Since $\CE_s$ and $E_0$ have CM by the same CM-field there exists an isogeny $\phi:\CE_s\rightarrow E_0$. We write $\phi_{dR}$ and $\phi_{\crys}$ for the morphisms induced from $\phi$ in the respective cohomology theories. 

Functoriality in the de Rham-crystalline comparison implies that \begin{equation}\label{eq:isogcompat}\CP_v(s)\circ \phi_{dR}=\phi_{\crys}\circ\CP_v(0).\end{equation} Note also that $\phi_{\crys}$ is induced from the reduction $\tilde{\phi}_v\in \End(\tilde{E}_{0,v})$ of $\phi$. Thus by \Cref{lemmacrysordbasis} we get that there exist $\zeta_v$, $\xi_v$ such that $[\phi_{\crys}]=\begin{pmatrix}	\zeta_v&0\\0&\xi_v\end{pmatrix}$ for the matrix of $\phi_{\crys}$ with respect to the basis of $H^1_v(E_0)$ fixed above.

Writing $\phi_{dR}\omega_0=a\cdot \omega_s$ and $\phi_{dR}\eta_0=c\cdot \omega_s+d\cdot \eta_s$ we may thus translate \eqref{eq:isogcompat} into the following equation between period matrices
\begin{equation}\label{eq:isogord1}	\begin{pmatrix}
			\iota_v(a)&0\\ \iota_v(c)&\iota_v(d)
	\end{pmatrix}\CP_v(s)=\CP_v(0)\begin{pmatrix}\zeta_v&0\\0&\xi_v\end{pmatrix}\end{equation}

Setting $y_{i,j}:=\iota_v(Y_{i,j}(x(s)))$, since $\CP_v(s)= (y_{i,j})\cdot \CP_v(0)$, we may rewrite the above as 
\begin{equation}
	\begin{pmatrix}
		\iota_v(a)&0\\ \iota_v(c)&\iota_v(d)
	\end{pmatrix}\cdot (y_{i,j})= \CP_v(0)\cdot \begin{pmatrix}\zeta_v&0\\0&\xi_v\end{pmatrix}\cdot \CP_v(0)^{-1}.
\end{equation}Since $\CP_v(0)=\diag(\varpi_v,\varpi_v^{-1})$ is diagonal, we easily get 
\begin{equation}
	\begin{pmatrix}
		\iota_v(a)&0\\\iota_v({c})&\iota_v(d)
	\end{pmatrix}\cdot (y_{i,j})=\begin{pmatrix}
	\zeta_v&0\\0&\xi_v
	\end{pmatrix}.
\end{equation}From this, given that $a\neq 0$ since $\phi_{dR}$ is invertible, we find $y_{1,2}=0$. \end{proof}

\begin{remark}\label{remarkordinaryproximity}
	Let $s$ be such that $\CE_s$ is CM. We note that $s$ will be $v$-adically close to $s_0$ with respect to some $v$ that is ordinary for $E_0$ if and only if the elliptic curves $E_0$ and $\CE_s$ have the same CM-field as their algebra of endomorphisms. Indeed $\End_{\bar{\Q}}(\CE_P)\hookrightarrow \End_{\bar{\mathbb{F}}_{p(v)}}^{0}(\tilde{E}_{0,v})$, for $P\in\{s,s_0\}$, since $\tilde{E}_{0,v} =\tilde{E}_{s,v}$ by our conventions on proximity. On the other hand, since $\CE_P$ is CM and $\tilde{E}_{0,v}$ is ordinary these algebras have to be equal for dimension reasons.
	
	In other words, $v$-adic relations among the values of our family of G-functions at points of interest for places $v$ that are ordinary for $E_0$ are only pertinent in the context described in \Cref{thmordrels}. 
\end{remark}
\subsection{Relations at supersingular places}\label{section:cmsuperpadic}

We retain the notation and conventions of \Cref{section:finrelsnotation} for our family of elliptic curves $f:\CE\rightarrow S$ and point $s_0\in S(K)$ whose fiber $E_0:=\CE_{s_0}$ is a CM elliptic curve. We let $F_0:=\End_{\bar{\Q}}^{0}(E_0)$ and consider the set 
\begin{center}$\Sigma_{\ssing,\unr}(s_0):=\{w\in\Sigma_{K,f}:w \text{ is supersingular for }E_0\text{ and } v\not| \disc(F_0) \}$,\end{center}of supersingular places of $E_0$ that are over primes unramified in $F_0$.


Here we restate a theorem of Y. Andr\'e which is the natural analogue of \Cref{thmordrels} for places of supersingular reduction of the ``center'' $E_0$. 
\begin{prop}[Y. Andr\'e \cite{andremots}, \cite{andremsj}]\label{propandresupercm}
	Let $s\in S(\bar{\Q})$ be such that $\CE_s$ is also a CM elliptic curve and consider the set of places \begin{center}
		 $\CP(s,s_0):=\{v\in \Sigma_{K(s),f}:\exists w\in\Sigma_{\ssing,\unr}(s_0), v|w \text{ and } s \text{ is }v\text{-adically close to }s_0\}\}$.
\end{center}
	
	Then there exists a polynomial $R_{s,\ssing,\unr}\in \bar{\Q}[X_{i,j}:1\leq i,j\leq 2]$ and an effectively computable positive constant $c_{\ssing}$ depending at most on $f:\CE\rightarrow S$ and $s_0$, such that the following hold
	\begin{enumerate}
		\item $R_{s,\ssing,\unr}$ has coefficients in the compositum $L_s$ of $K(s)$ with $\End_{\bar{\Q}}^{0}(\CE_s)$,
		
		\item $\iota_v(R_{s,\ssing,\unr}(Y_G(x(s))))=0$ for all $v\in\Sigma_{L_s,f}$ for which there exists $w\in\CP(s,s_0)$ with $v|w$, 
		
		\item $R_{s,\ssing,\unr}$ is homogeneous with $\deg(R_{s,\ssing,\unr})\leq c_{\ssing}\cdot|\CP(s,s_0)|$, and 
		
		\item $R_{s,\ssing,\unr}\notin \langle X_{1,1}X_{2,2}-X_{1,2}X_{2,1}-1 \rangle $.
	\end{enumerate}\end{prop}
\begin{proof}
	We give a quick sketch of a proof here, especially since the same ideas appear in greater technical detail in the proof of the corresponding fact for QM abelian surfaces in the sequel \cite{papaspadicpart3}. 
	
	Let $d_s$ be such that $\End_{\bar{\Q}}^{0}(\CE_s)=\Q(\sqrt{-d_s})$. Let us also fix $v\in \Sigma_{L_s,f}$ such that $s$ is $v$-adically close to $s_0$ and $v$ is supersingular for $E_0$ and unramified in $\End_{\bar{\Q}}^{0}(E_0)=:\Q(\sqrt{-d_0})$. In $\S 5$ of \cite{andremots} Andr\'e finds $a_s\in L_s$ (denoted in $\sigma$ in loc. cit.) and $m_{s,v}\in L_s$ (denoted $m_v$ in loc. cit. ) such that the polynomial defined by \begin{equation}\label{eq:andresupersingularpoly}
		R_{s,v}:=(m_{s,v}+2\sqrt{d_sd_0})X_{1,1}(X_{2,2}+a_sX_{{1,2}})-(m_{s,v}-2\sqrt{d_sd_0})X_{1,2}(X_{2,1}+a_sX_{1,1}),
	\end{equation}satisfies $\iota_v(R_{s,v}(Y_G(x(s))))=0$.
	
	Consider the matrices $S(n,l):=\begin{pmatrix}n&l\\0&\frac{1}{n}\end{pmatrix}\in \SL_2(\C)$ and  $T(n):=\begin{pmatrix}0&\frac{1}{n}\\-n&0\end{pmatrix}\in \SL_2(\C)$. Assume that we had $R_{s,v}\in \langle X_{1,1}X_{2,2}-X_{1,2}X_{2,1}-1 \rangle$. Then $R_{s,v}(S(n,l))=0$ implies $nl(4a_s\sqrt{d_0d_s})+(m_{s,v}+2\sqrt{d_0d_s})=0$ while $R_{s,v}(T(n))=0$ gives $m_{s,v}-2\sqrt{d_0d_s}=0$. Since the first of these holds for all $n$, $l\in \C$ we get $a_s=0$, $m_{s,v}+2\sqrt{d_0d_s}=0$, and  $m_{s,v}-2\sqrt{d_0d_s}=0$. The last two give a contradiction since $d_sd_0\neq 0$. 
	
	The polynomial we want is nothing but $R_{s,\ssing,\unr}:=\underset{v}{\Pi} R_{s,v}$, with $R_{s,v}$ as in \eqref{eq:andresupersingularpoly}, where the product ranges over all the $v\in\Sigma_{L_s}$ which are such that $|x(s)|_v<\min\{1,R_v(Y_G)\}$ and $v|w$ for some $w\in \Sigma_{\ssing,\unr}(s_0)$. All properties follow trivially by definition.\end{proof}

\subsection{Relations at archimedean places}\label{section:cmarchimedean}

In the archimedean case the analogous statement to \Cref{propandresupercm}, due to F. Beukers, is the following:
\begin{prop}[Theorem $3.3$, \cite{beukers}]\label{propbeukersarch}
		Let $s\in S(\bar{\Q})$ be such that $\CE_s$ is also a CM elliptic curve. Then there exists a polynomial $R_{s,\arch}\in \bar{\Q}[X_{i,j}:1\leq i,j\leq 2]$ such that the following hold
	\begin{enumerate}
		\item $R_{s,\arch}$ has coefficients in the compositum $L_s$ of $K(s)$ with $\End_{\bar{\Q}}^{0}(\CE_s)$,
		
		\item $\iota_v(R_{s,\arch}(Y_G(x(s))))=0$ for all $v\in\Sigma_{L_s,\infty}$ for which $|x(s)|_v<\min\{1,R_v(Y_G)\}$,
		
		\item $R_{s,\arch}$ is homogeneous with $\deg(R_{s,\arch})\leq c_{\arch}\cdot [\Q(s):\Q]$, where $c_{\arch}$ is some absolute constant, and 
		
		\item $R_{s,\arch}\notin \langle X_{1,1}X_{2,2}-X_{1,2}X_{2,1}-1 \rangle $.
		\end{enumerate}\end{prop}
	
	\begin{proof}Our reformulation of Beukers' result follows practically identically from the same argument as above. Given $v\in\Sigma_{L_s,\arch}$ with $|x(s)|_v<\min\{1,R_v(Y_G)\}$ Beukers shows that there exist $a_s$, $m_{s,v}$, $n_{s,v}\in L_s$ such that for the polynomials 	\begin{center}
					$ R_{s,v,1}:=X_{1,1} (X_{2,2}+a_sX_{1,2})-m_{s,v}(X_{1,1}X_{2,2}-X_{1,2}X_{2,1})$, and 
		
				$ R_{s,v,2}:=X_{1,2} (X_{2,1}+a_sX_{1,1})-n_{s,v}(X_{1,1}X_{2,2}-X_{1,2}X_{2,1})$\end{center}
we have $\iota_v(R_{s,v,j}(Y_G(x(s))))=0$.

First note that, since $\det(Y_G(x))=1$ generically, if $R_{s,v,1}\in  \langle X_{1,1}X_{2,2}-X_{1,2}X_{2,1}-1 \rangle $ we would have, arguing as in the previous proof, $a_{s}=0$ and $m_{s,v}=1$. This may be seen by ``testing'' at the family $S(n,l)$ for example. But then we get that $X_{1,1}X_{2,2}-1\in \langle X_{1,1}X_{2,2}-X_{1,2}X_{2,1}-1 \rangle $ which is trivially false.

The polynomial we want is nothing but the product $\R_{s,\arch}:=\underset{v}{\Pi}R_{s,v,1}$, where the $v$ range over those archimedean places of $L_s$ for which $|x(s)|_v<\min\{1,R_v(Y_G)\}$.\end{proof}
\subsection{Relations at ramified places}\label{section:ramifiedprimescm}

The missing piece in completing the picture in \Cref{naiverelations} is the construction of relations at points $s$ with $\CE_s$ CM elliptic curves among the values of our G-functions at places $v$ that are over primes $p$ ramified in $F_0$. 

\begin{prop}\label{propramifiedcm}
	Let $f:\CE\rightarrow S$ be as in \Cref{section:finrelsnotation}. Let $s\in S(\bar{\Q})$ be such that $\CE_s$ is CM.
	
	Then, there exists a polynomial  $R_{s,\ram}\in \bar{\Q}[X_{i,j}:1\leq i,j\leq 2]$ such that the following hold
	\begin{enumerate}
		\item $R_{s,\ram}$ has coefficients in some $L_s/K(s)$ with $[L_s:K(s)]$ bounded by an absolute constant,
		
		\item $\iota_v(R_{s,\ram}(Y_G(x(s))))=0$ for all $v\in\Sigma_{L_s,f}$ for which $|x(s)|_v<\min\{1,R_v(Y_G)\}$ and $v|\disc(F_0)$,
		
		\item $R_{s,\ram}$ is homogeneous with $\deg(R_{s,\ram})\leq c_{\ram}[K(s):\Q]$, where $c_{\ram}$ is some absolute constant, and 
		
		\item $R_{s,\ram}\notin \langle X_{1,1}X_{2,2}-X_{1,2}X_{2,1}-1 \rangle $.
\end{enumerate}\end{prop}
\begin{proof}We begin by noting that in any case here we must have that $F_s=F_0$, i.e. $\CE_s$ and $E_0$ have CM by the same CM field. To see this we start by writing $d_P:=\disc(\End_{\bar{\Q}}(\CE_P))$ for $P\in\{s,s_0\}$. If $d_s=d_{s_0}$ then we are done. If $d_s\neq d_{s_0}$ our assertion follows from \Cref{lemproximityramified}.
	
Now we may argue much as  in the proof of \Cref{thmordrels} to find an isogeny $\phi:\CE_s\rightarrow E_0$ defined over $\bar{\Q}$. The extension $L_s$ will be the compositum of $K(s)$ with the field of definition of this $\phi$. The degree $[L_s:K(s)]$ may be bounded by an absolute constant, see for example \cite{silverberg}. From now on we fix $v\in \Sigma_{L_s,f}$ which is such that $v|\disc(F_0)$ and such that our $s$ is $v$-adically close to $s_0$.

We write $\phi_{dR}$ and $\phi_{\crys}$ for the pullbacks of $\phi$ on the level of de Rham and Crystalline cohomology respectively. Once again we will have that $\phi_{\crys}$ will be the pullback on the level of Crystalline cohomology of the reduction modulo $v$ of $\phi$, which we denote by $\tilde{\phi}_v$. As in the proof of \Cref{thmordrels}, by our conventions on $v$-adic proximity in \Cref{defnvadicprox}, we get that $\tilde{\phi}_v\in \End(\tilde{E}_{0,v})$ since $\tilde{\CE}_{s,v}=\tilde{E}_{0,v}$.

Once again we get \eqref{eq:isogcompat}, which we may rewrite, much as in \eqref{eq:isogord1}, as 
\begin{equation}\label{eq:isogram1}	[\phi]_{dR}\CP_v(s)=\CP_v(0)[\phi]_{\crys},
\end{equation}where $[\phi]_{*}$ denotes the matrix induced from $\phi_{*}$, for $*\in \{dR,\crys\}$, with the implied bases on each side of the de Rham-Crystalline comparison isomorphism.
	
Working much as in the proof of \Cref{thmordrels}, we may assume that $[\phi]_{dR}=\begin{pmatrix}\iota_v(a)&0\\ \iota_v(c)&\iota_v(d)\end{pmatrix}$. Setting $y_{i,j}:=\iota_v(Y_{i,j}(x(s)))$, since $\CP_v(s)= (y_{i,j})\cdot \CP_v(0)$, we may rewrite \eqref{eq:isogram1} as 
	\begin{equation}\label{eq:ramif2}
		\begin{pmatrix}
			\iota_v(a)&0\\ \iota_v(c)&\iota_v(d)
		\end{pmatrix}\cdot (y_{i,j})= \CP_v(0)\cdot [{\phi}]_{\crys} \cdot \CP_v(0)^{-1}.
	\end{equation}
	
Taking traces in \eqref{eq:ramif2} we conclude that 
	\begin{equation}\label{eq:ramif3}
		\iota_v(a Y_{1,1}(x(s))+c Y_{1,2}(x(s))+dY_{2,2}(x(s)))=t_v:=\tr([\phi]_{\crys}).
	\end{equation}
	
Practically, all we need at this point is to note that $t_v\in \Q$. This follows from the fact that $\phi_{\crys}$ is the pullback of the endomorphism $\tilde{\phi}_v$ of $\tilde{E}_{0,v}$. See the proof of Lemma $3.5$ in \cite{daworrpap} for more details.

At this point, much as in the proof of \Cref{propandresupercm}, we set $R_{s,v}:=(aX_{1,1}+c Y_{1,2}+dY_{2,2})^2-t_v\det(Y_{i,j})$. Note that this way $R_{s,v}$ is a homogeneous degree $2$ polynomial and from \eqref{eq:ramif3} we have $\iota_v(R_{s,v}(Y_G(x(s))))=0$, since $\det(Y_{i,j}(x))=1$ on the level of power series. The polynomial we want is nothing but $R_{s,\ram}:=\prod R_{s,v}$, where the product is over all $v\in\Sigma_{L_s,f}$ that divide $\disc(F_0)$ and are such that $s$ is $v$-adically close to $s_0$.

To check that $R_{s,\ram}$ defines a non-trivial relation among the values of $Y_G$ at $s$ we may work again exactly as in \Cref{propandresupercm}. In more detail, we need only check that $R_{s,v}$ is not in the ideal $\langle X_{1,1}X_{2,2}-X_{1,2}X_{2,1}-1\rangle$. Testing with the same $S(n,l)$ as in \Cref{propandresupercm} is easily seen to force $a=0$ which contradicts the invertibility of the matrix $[\phi]_{dR}$ corresponding to our isogeny $\phi$.\end{proof}

\begin{remark}
	We note that we need no particular property of the basis of the Crystalline cohomology of the central fiber $H^1_v(E_0)$ in this construction. Rather everything follows from the existence of the isogeny $\phi$, which is guaranteed to us by \Cref{lemproximityramified}.
\end{remark}
%\part{Applications}
\section{Height bounds and applications}\label{section:applications}

In this final section we apply Andr\'e-Bombieri's ``Hasse principle for the values of G-functions'' to reach the height bounds alluded to in \Cref{section:heightboundsintro}. In the process, we attempt to isolate the difficulty of obtaining the ``ideal height bounds'' one would want here to conclude an effective version of \cite{siegel}.

\subsection{Height bounds}\label{cmheightboundproof}

We start by presenting a height bound for the singular moduli of CM elliptic curves whose CM is given by the same CM field as that of a fixed CM elliptic curve. This takes the following form:
\begin{prop}\label{propisogenouscm}
	Let $j_0$ be the $j$-invariant of a CM-elliptic curve $E_0$ defined over a number field $K$ and consider the set \begin{center}
		$A(j_0):=\{j: \End^{0}_{\bar{\Q}}(E_j)=\End^{0}_{\bar{\Q}}(E_0)\}$,
	\end{center}of $j$-invariants whose corresponding elliptic curves, i.e. $E_j$, have CM by the same CM field as $E_0$.
	
	Then, there exist effectively computable constants $C_1$, $C_2>0$, depending only on $K$ and $j_0$, such that
	\begin{center}
		$h(j)\leq C_1 +C_2\log ([\Q(j):\Q])$, 
	\end{center}for all $j\in A(j_0)$.\end{prop}
\begin{proof} Without loss of generality, by possibly replacing $K$ by a finite extension $K_0/K$ whose degree depends on $j_0$, we may assume that $E_0$ has everywhere good reduction over $K$, see \cite{serretate}.
	
	Let $K(j)$ be the compositum of $K$ and $\Q(j)$. By \cite{silverberg} there exists an extension $K'(j)/K(j)$, with $[K'(j):K(j)]$ bounded by an absolute constant, such that $\End^{0}_{\bar{\Q}}(E_j)=\End^{0}_{K'(j)}(E_j)$ and $\End^{0}_{\bar{\Q}}(E_0)=\End^{0}_{K'(j)}(E_0)$. Since $\End^{0}_{\bar{\Q}}(E_j)=\End^{0}_{\bar{\Q}}(E_0)$ the two elliptic cures are isogenous over $\bar{\Q}$. By Lemma $6.1$ of \cite{masserwuisogellcurves} there exists an extension $L(j)/K^{\prime}(j)$ with $[L(j):K^{\prime}(j)]\leq 12$ such that such an isogeny $\phi$ is defined over $L(j)$.
	
	By Theorem $1.1$ of \cite{pazuki} there exist effectively computable positive constants $c_1$ and $c_2$, depending at most on $j_0$, such that 
	\begin{equation}h(j)\leq c_1+c_2\log(\deg(\phi)).\end{equation}
	On the other hand, by Theorem $1.4$ of \cite{gaudronremondpolarisations} there exists such an isogeny with $\deg(\phi)\leq \kappa(E_0\times_{K}L(j))$ where $\kappa(A)$ is the quantity defined at the beginning of the introduction of loc. cit.. By definition here we have \begin{center}
		$\log(\kappa(E_0\times_KL(j)))\leq c_1^{\prime}+c_2^{\prime}\log([L(j):\Q])+$
		
		$+c_3^{\prime}\log(\max\{h_F(E_0\times_KL(j)),\log([L(j):\Q]),1\})$.
	\end{center}
	
	Now trivially, $\log(\max\{h_F(E_0\times_KL(j)),\log([L(j):\Q]),1\})\leq \max\{h_F(E_0\times_KL(j)),\log([L(j):\Q]),1\}\leq h_F(E_0\times_KL(j))+1+\log([L(j):\Q])$. Since $E_0$ has, by assumption, everywhere good reduction, and thus also semi-stable reduction, over $K$, we have by basic properties of the stable Faltings height that $h_F(E_0\times_KL(j))=h_F(E_0)$. The conclusion now follows trivially from the above.\end{proof}
	
	\begin{remark}
		We believe that the above is known to experts. We have opted to include this proof here for two reasons. 
		
		Firstly, the tools used in the proof are either standard in proofs of lower bounds on Galois orbits, or, in the case of Pazuki's result, follow much that same ``mantra'' as the output of the G-functions method. In other words, the ``mantra'' of bounding the Weil height of a point of interest by arithmetic data associated to our geometric objects.
		
		The other reason we have opted to include this here, is that it isolates further the problem of obtaining the height bounds envisioned by Y. Andr\'e in connection to the ``Effective Brauer-Siegel'' problem.
\end{remark}

Before stating our height bound we introduce some notation.
\begin{definition}\label{defnproblematic}
	Let $j_0\neq 0$ be a fixed singular modulus whose corresponding elliptic curve has CM by the CM field $F_0$. Given another singular modulus $j$ whose corresponding CM elliptic curve has CM given by a CM field $F_j\neq F_0$ we define
	 \begin{equation}\CP(j,j_0):=\{ v\in \Sigma_{\Q(j,j_0),f}: v(j-j_0)>0, \tilde{E}_{0,v}\text{ is supersingular}\}.
		\end{equation}We call $\CP(j)$ the set of \textbf{supersingular places of proximity} of $j$ to $j_0$.
	\end{definition}

The main output of the G-functions method in this setting is the following height bound:
\begin{theorem}\label{heightboundcm}Let $j_0:=54000$ and let $\CP(j):=\CP(j,54000)$. Then, there exist effectively computable positive constants $c_1$ and $c_2$, such that
\begin{equation}\label{eq:htboundmain}h(j)\leq c_1 ([\Q(j):\Q] +|\CP(j)| )^{c_2},\end{equation}
for all singular moduli $j$ whose corresponding elliptic curve has CM by a field other than $\Q(\sqrt{-3})$.\end{theorem}

\begin{proof} We let $f:\CE\rightarrow S$, where $S:=\mathbb{A}^1\backslash \{0,1728\}$, be the $j$-family of elliptic curves, see for example page $47$ in \cite{silvermanell}. We also fix $j_0=s_0=54000\in S(\Q)$, which is such that\footnote{See \cite[\href{https://www.lmfdb.org/EllipticCurve/2.0.3.1/144.1/CMa/2}{Elliptic Curve 144.1-CMa2}]{lmfdb}.} $E_0:=\CE_{s_0}$ is a CM elliptic curve with CM field $\Q(\sqrt{-3})$. To the pair $(f,s_0)$ we associate the local parameter $x:=j-j_0$. To this picture, as discussed in \Cref{section:backgroundgfuns}, we may associate a family of G-functions. In this case, this will be nothing but a single matrix which we denote from now on by $Y_G$.
	
Throughout this proof let $j=s\in S(\bar{\Q})$ be another point for which $\CE_s$ is CM, with $\End^{0}_{\bar{\Q}}(\CE_s)\neq \Q(\sqrt{-3})$. We consider the set $\Sigma(s):=\{v\in \Sigma_{L_s}:|x(s)|_v<\min\{1,R_v(Y_G)\}\}$ of places of the extension $L_s$ considered in \Cref{propandresupercm} with respect to which $s$ is close to $s_0$. If $\Sigma(s)=\emptyset$ we conclude as in the proof of the main height bound of \cite{papasbigboi}. From now on we assume $\Sigma(s)\neq \emptyset$ and consider its subsets $\Sigma(s)_{\arch}:=\{v\in \Sigma(s):v|\infty\}$, $\Sigma(s)_{\ord}:=\{v\in\Sigma(s):v \text{ is ordinary for }E_{0} \}$, and $\Sigma(s)_{\ssing}:=\{v\in\Sigma(s):v  \text{ is supersingular for }E_0\}$.

Since $\End^{0}_{\bar{\Q}}(\CE_s)\neq \Q(\sqrt{-3})$ we can easily see that $\Sigma(s)_{\ord}=\emptyset$, see \Cref{remarkordinaryproximity}. Similarly, from \Cref{lemproximityramified} we know that if $v\in \Sigma(s)_{\ssing}$ and $p$ is the unique prime in $\Q$ which is divisible by $v$, then $p$ is unramified in $\Q(\sqrt{-3})$. In other words, $p\neq 3$ the unique prime ramified in $\Q(\sqrt{-3})$. In particular, we may use \Cref{propandresupercm} to get a polynomial $R_{s,\ssing}\in L_s[X_{i,j}:1\leq i,j\leq 2]$ such that $\iota_v(R_{s,\ssing}(Y_G(x(s))))=0$ for all $v\in\Sigma(s)_{\ssing}$. 

The places of the set $\Sigma(s)_{\arch}$ are ``dealt with'' by the polynomial $R_{s,\arch}$ of \Cref{propbeukersarch}. The polynomial $R_{s}:=R_{s,\arch}\cdot R_{s,\ssing}$ will satisfy \begin{center}
	$\iota_v(R_{s}(Y_G(x(s))))=0$
\end{center} for all places $v\in\Sigma(s)$, thanks to our comments above. This polynomial will correspond to a global non-trivial relation among the values of the G-functions of our family at the point $s$. We may now apply Andr\'e-Bombieri's ``Hasse principle for values of G-functions'', see Chapter $VII.5$ in \cite{andre1989g}. The result now follows since $\deg(R_s)$ is bounded by the quantity on the right of \eqref{eq:htboundmain}.\end{proof}

	We note that the choice of $j_0$ here is completely arbitrary, other than the condition $j_0\neq 0,1728$, as well as the choice of the CM field $\Q(\sqrt{-3})$. A different choice of $j_0$ would alter the constant $c_1$ in \eqref{eq:htboundmain}, since $c_1$ depends on arithmetic quantities of the fiber $\CE_{s_0}$.
	
The same exact argument as above would give us the slightly more general:
\begin{theorem}\label{htboundgfunsgeneral}
	Let $f:\CE\rightarrow S$ be a $1$-parameter family of elliptic curves defined over a number field $K$ and $s_0\in S(K)$ a point whose fiber $\CE_{s_0}$ has CM by the CM field $F_0$. Assume that the induced map $S\rightarrow Y(1)$ is not-isotrivial. 
	
	Then there exist effectively computable positive constants $C_1$, $C_2$, depending on $f:\CE\rightarrow S$ and $s_0$, such that 	
	\begin{equation}
		h(s)\leq C_1 ([K(s):K]+|\CP(j(s),j(s_0))|)^{C_2},
	\end{equation}for all $s\in S(\bar{\Q})$ whose fibers have CM given by some CM field other than $F_0$, and where $j(P)$ denotes the $j$-invariant of the fiber $\CE_P$.
\end{theorem}

\subsection{The effective Brauer-Siegel problem}\label{section:brauersiegel}

 One of Andr\'e's motivations towards his seminal work in \cite{andre1989g}, seems to have been, see in particular Remark $3$ on page $201$ of loc. cit., the search for an effective version of the following classical, and famously ineffective, result of C. L. Siegel:
\begin{theorem}[Siegel,\cite{siegelog}]\label{siegelogthm}Given $\epsilon>0$ there exists a constant $c(\epsilon)>0$ such that 
	\begin{equation}\label{eq:siegeleq}
		h(K)\geq c(\epsilon) |\disc(K)|^{\frac{1}{2}-\epsilon},
	\end{equation}for all imaginary quadratic number fields $K$.
\end{theorem}

Via standard results about CM elliptic curves, \Cref{siegelogthm} may be translated to a statement about CM elliptic curves. Given a singular modulus $j$ we let $E_j$ be the corresponding CM elliptic curve and set $D(j):=\disc(\End_{\bar{\Q}}(E_j))$ to be the discriminant of the ring of $\Q$-endomorphisms of $E_j$. An effective answer to \Cref{siegelogthm} would be equivalent\footnote{We point the interested reader to the introduction of \cite{papaseffbrsieg} for more on this.} to finding an effective constant $c(\epsilon)>0$ such that 
\begin{equation}\label{eq:siegelrestated}
	[\Q(j):\Q]\geq c(\epsilon) |D(j)|^{\frac{1}{2}-\epsilon}.
\end{equation}

The proof of Theorem $1.5$ in \cite{papaseffbrsieg}, replacing the usage of the height bounds in Theorem $2.15$ there by the height bound in \Cref{heightboundcm}, leads to the existence of effectively computable constants $c_4$, $c_5$ for which 
\begin{equation}\label{eq:brauersiegel1}
	|D(j)|\leq c_4\max\{[\Q(j):\Q],[\Q(j):\Q]+|\CP(j)|\}^{c_5},
\end{equation}for all singular moduli $j$ for which $\End^{0}_{\bar{\Q}}(E_j)\neq \Q(\sqrt{-3})$. Compare, in particular, equation $(52)$ in loc. cit. with \eqref{eq:brauersiegel1}. Furthermore note here that, were we to use \Cref{propisogenouscm} instead of \Cref{heightboundcm}, we would get effective constants dealing with the $j$ with $\End^{0}_{\bar{\Q}}(E_j)=\Q(\sqrt{-3})$, albeit with a fractional power on the right of \eqref{eq:siegelrestated} slightly worse than $\frac{1}{2}$.

The obvious question that now arises is if the contribution of $|\CP(j)|$ can be made ``negligible''. What one would want here is to find effectively computable constants $C_{1}$, $C_2$, and $C_3>0$ such that, for all singular moduli $j$ with $\End^{0}_{\bar{\Q}}(E_j)\neq \Q(\sqrt{-3})$ we have
	\begin{equation}\label{eq:conjecturalbound}
		|\CP(j)|\leq C_{1}[\Q(j):\Q]^{C_2}\cdot D(j)^{\frac{1}{C_3}},\end{equation}with $C_3>c_5$.


We close off with some remarks on the constant $c_5$ of \eqref{eq:brauersiegel1}. First of all, in this case, the constant $c_5$ will be the product of the constant $c_1$ outputted by Andr\'e-Bombieri's ``Hasse-Principle'' and a constant $c_1'$ outputted by Masser-W\"ustholz's ``Endomorphism estimates'', see \cite{mwendoesti}.

In \cite{mwendoesti} Masser and W\"ustholz note that the optimal value of the constant $c_1'$ is of the magnitude $2+\epsilon$ in the case of elliptic curves. Andr\'e's method on the other hand would give $c_1=9$. We note that, following Remark $2$ on page $138$ of \cite{andre1989g}, we may even take $c_1=4$, i.e. using the height bound for ``strongly non-trivial relations'' in the vocabulary of Andr\'e's version of the ``Hasse principle'' in loc. cit.. All in all, it would therefore suffice for our purposes to have a version of \eqref{eq:conjecturalbound} with $C_3>8$.

\begin{remark}The above, perhaps naive, sketch of a strategy towards an effective version of Siegel's result, is at first glance at least, crude from the perspective of the ``G-functions method''. More precisely, the problematic quantity in the height bounds in \Cref{heightboundcm} is there to control the way that the relations outputted by \Cref{propandresupercm} depend on the place in question. A, perhaps, more reasonable approach would be to try to directly bound the number of the possible polynomials that might appear in \Cref{propandresupercm}, by, say, the same quantity as in \eqref{eq:conjecturalbound}.

This approach was, in principle at least, considered by Beukers, see in particular Theorem $5.4$ in \cite{beukers}. Using Beukers' results, the height bound in \Cref{heightboundcm} would look like \begin{equation}\label{eq:badheightbound}
	h(j)\leq c_0 [\Q(j):\Q]^4\cdot D(j)^2.
\end{equation}Were we to pair this with Masser-W\"ustholz's estimates, we would reach a clear impasse. \end{remark}

\appendix	

\bibliographystyle{alpha}
\bibliography{biblio}
\end{document}
